\section{Introduction}\label{sec:introduction}

Viterbo's conjecture asserts that among convex bodies of fixed volume the
ball has the largest symplectic capacity~\cite{Viterbo2000}. We work in
$\mathbb R^4$ with the Ekeland--Hofer--Zehnder capacity $c_{\rm EHZ}$: for a
convex body $K$ it is the least action of a closed characteristic on
$\partial K$, where for polytopes the closed characteristics are replaced by
generalized characteristics. The \emph{systolic ratio}
\[
 \operatorname{sys}(K)=\frac{c_{\rm EHZ}(K)^2}{2\operatorname{vol}_4(K)}
\]
equals one for the ball, and the conjecture asserts
$\operatorname{sys}(K)\leq1$ for every convex body $K\subset\mathbb R^4$.
Haim-Kislev and Ostrover disproved it with a Lagrangian
product~\cite{HaimKislevOstrover2024}. For planar convex bodies $K_q$ and
$K_p$, the Lagrangian product $K_q\times_L K_p$ places $K_q$ in the
$(q_1,q_2)$-plane and $K_p$ in the $(p_1,p_2)$-plane. With $P$ a regular
pentagon and $R(\theta)$ the rotation by $\theta$, the HKO body
$K_{\rm HKO}=P\times_L R(-\pi/2)P$ has systolic ratio
$(3+\sqrt5)/5\approx1.047$.

This thesis fixes a toolbox and tests how far it gets on questions around
Viterbo's conjecture after this counterexample. The toolbox consists of
high-performance computation on Augsburg's LICCA cluster, proof by
computation in exact arithmetic, standard data-science methods, and AI
agents. The agents ran the data-science experiments and did large parts of
the theory work; Section~\ref{sec:use-of-ai} describes how. The questions
were not chosen first: they came out of the work.

Standard data science found no new body with systolic ratio above one.
In random samples of generic polytopes and of Lagrangian products of
polygons, and in local searches started from generic ten-facet polytopes,
every computed systolic ratio is below one. Perturbations of $K_{\rm HKO}$
have systolic ratio above one, but they are small variations of the known
counterexample. The data did show a strong inverse association between
systolic ratio and symplectic ridge area, the total symplectic area of the
two-dimensional faces normalized by volume. In knot theory, patterns found
in computed data have suggested results that then inspired theory; here AI
agents both look for the patterns and do much of the theory.

Randomly sampled perturbations of $K_{\rm HKO}$ never increased its
systolic ratio. This raised the first question that the thesis answers by
a theorem: is $K_{\rm HKO}$ a local maximum of the systolic ratio? We prove
that it is among polytopes with exactly ten facets, and that the same space
contains a second local maximum, a hexagon--quadrilateral product with
systolic ratio one. Other results answer questions that arose the same way:
an explicit formula for the systolic ratio of $P\times_L R(\theta)P$ as a
function of $\theta$, the sharp bound $(3+\sqrt5)/5$ for products of
independently linearly deformed regular pentagons, and a capacity-increasing rotation of any product of
planar convex bodies out of symplectic position. The capacity computations
rest on Haim-Kislev's finite formula~\cite{HK2017}, a reduction for
Lagrangian products, and a second algorithm built on the local geometry of
Chaidez and Hutchings~\cite{CH2021}.

The subsections below first recall how capacities connect embeddings,
dynamics and convex geometry, then introduce the local theorem, the finite
capacity calculations and the data-science results, and end with suggested
reading orders.
